% sup3_04_S14b: Appendix on Ba3 Prop. 1, part 3, step S14, second half: the expanded expression of a class, and its decomposition into renormalised classes.

\subsubsection{S14, II: the expanded expression and its classes}\label{sup3:ssec:S14b}

\begin{lemma}[the expanded expression]\label{sup3:lem:S143}
Let $\mathcal C$ be a connected component of a class after the cube split, with all scalar lines $G_k(\square,\tilde B)$ (P-iii) and all old vertices localised in the cube $Q_0\subset\square_1$ of \cref{sup3:cor:S114}(c). Expand every scalar line by \cref{sup3:lem:S142}(b$'$) with fixed $n_1,n_2\ge2$ (the orders for which the last term of (b$'$) is a regular kernel of both variables, as in the fourth part), so with at most $N_B=n_1+n_2-1$ insertions in a main term, and every old vertex by \cref{sup3:lem:S141}: its dependence on $\tilde B$ is through $D^\eta_{\tilde B}$ in \cite[(1.8)--(1.9)]{Ba3} and through $U(\tilde B(\Gamma))$ in \cite[(1.13)--(1.15)]{Ba3}, while \cite[(1.6), (1.7), (1.10), (1.11)]{Ba3} and the counterterm vertices do not depend on $\tilde B$. Then
\[
E(\mathcal C)=\sum_{G\in\mathcal C}\ \sum_{\tau\in T(G)}E(G_\tau),
\]
a finite sum, where a \emph{pattern} $\tau$ chooses, for each scalar line, one main term of (b$'$) together with one vertex operator of \cref{sup3:lem:S142}(c) per insertion, or the remainder term, and for each old vertex one term of its expansion. Each $G_\tau$ has:
\begin{enumerate}[label=\textup{(\roman*)},nosep,leftmargin=*]
\item scalar lines $G_k(\square,B_0)$ and vector lines unchanged (they do not depend on $\tilde B$); the external fields unchanged; each line with the remainder choice replaced by a pair of external scalar legs carrying the two-variable field $(G_0V_k)^{n_1}G_k(\square,\tilde B)(V_kG_0)^{n_2}$;
\item old vertices of the same type and the same localisation box $Q_v$, with $B_0$ in place of $\tilde B$, and extra $\tilde B'$-legs at the vertex's own bond or contour, inside $Q_v$, so that no relocalisation is needed. ``Same type'' is meant collectively, as \cite[p.~433]{Ba3} says it (``Vertices are now of the same type (1.6)--(1.11) and (1.13)--(1.15), but with $B_0$ instead of $\tilde B$ and with the new external vector field $\tilde B'$ beside the old one''): the term $F_{1,k}(\tilde B'_b)U(B_{0,b})\phi(b_+)\cdot q^m\phi(b_-)$ of an expanded \cite[(1.8)]{Ba3} or \cite[(1.9)]{Ba3} vertex has no derivative leg and $n'+j$ vector legs, i.e.\ the shape of \cite[(1.10)]{Ba3} or \cite[(1.11)]{Ba3}; its degree is the old one plus $j$;
\item new vertices, at most $N_B$ on each old scalar line, each of the form of \cref{sup3:lem:S142}(c): exactly two internal lines, both on the chain of that old line, no $A'$-leg, lying on bonds and blocks of $\square$ (their label boxes may leave $\square$ by at most $1$). An old scalar line may have one old end and one external scalar leg; its chain then has a single old end;
\item $|T(G)|\le\big[(N_B+2)(6(\bar n+1)^2)^{N_B}\big]^{L(\bar n)}\,(2(\bar n+1))^{2dV_{\mathrm{mer}}(\bar n)}$, with $L(\bar n)\le2V_{\mathrm{mer}}(\bar n)$ the maximal number of scalar lines: a bound in terms of $\bar n$, $N_B$, $d$ only;
\item every $\tilde B'$-factor is $e_k\tilde B'_b$ or $e_k\tilde B'(\Gamma)$, with $|e_k\tilde B'_b|\le e_kC_{13}p(t)$ and $|e_k\tilde B'(\Gamma)|\le e_kdC_{13}p(t)$ (\cref{sup3:lem:S13}); the $R$-factors have norm $\le1$;
\item $V_{\mathrm{mer}}(G_\tau)\le V_{\mathrm{mer}}(\bar n)+N_BL(\bar n)\le V_{\mathrm{mer}}(\bar n)(1+2N_B)$; the label count after the expansion is bounded by a function $V^\sharp_{\mathrm{lab}}(\bar n)$ of $\bar n$ (and the fixed $N_B$) only, since each new vertex carries at most $\bar n+1$ labels.
\end{enumerate}
\end{lemma}

\begin{proof}
$E(G)$ is multilinear in the factors that depend on $\tilde B$: each scalar line, and each $D^\eta_{\tilde B}$ or $U(\tilde B(\Gamma))$ at an old vertex. Replace each by its finite expansion (\cref{sup3:lem:S141,sup3:lem:S142}) and distribute; this is the sum over $\tau$. (i)--(iii) read off the terms: an inserted vertex operator of \cref{sup3:lem:S142}(c) lies on a bond or block of $\square$ (the operators act on fields of $\square$), its two scalar legs are joined to its neighbours on the chain, and its label boxes, of side $\le2$ around a point of $\square$, leave $\square$ by at most $1$. (iv) Per scalar line there are $n_1+n_2=N_B+1$ main terms and one remainder, and each insertion has at most $6(\bar n+1)^2$ vertex operators; per old vertex at most $2d$ bond or contour factors of $\tilde B$-type, each with at most $\bar n+2\le2(\bar n+1)$ terms; a graph has at most $4V_{\mathrm{mer}}/2$ scalar lines. (v) is \cref{sup3:lem:S13}(a), (c) with $|\Gamma|_\eta\le d$. (vi) Each line gains at most $N_B$ vertices.
\end{proof}

\paragraph{The regrouping into renormalised classes.} The class rules of \cite[Sect.~2]{Ba3} (\cref{sup2:ssec:S7}): divergent graphs ($D\le0$) are built of \cite[(1.6)--(1.8), (1.10)]{Ba3} only, and a graph with an $R$-vertex, a block vertex \cite[(1.13)--(1.15)]{Ba3}, a \cite[(1.7)]{Ba3} vertex with finite counterterm or more than four external legs has $D>0$ \cite[Cor.~2.3]{Ba3}; the two-scalar-leg and $\phi\phi A$ classes are ``the smallest sets of graphs, symmetric with respect to permutation of external scalar field legs and with each divergent subgraph renormalized'' \cite[p.~429]{Ba3}, two-scalar-leg classes with their mass counterterm \cite[p.~430]{Ba3}; the $2A$ classes contain ``all graphs connected with these graphs by Ward--Takahashi identities'', or, ``if we refrain from taking the smallest classes'', ``all possible connected graphs of the considered type \dots\ of a fixed order'' \cite[p.~431]{Ba3}; $1A$ and $3A$ graphs and differentiated two-scalar-leg graphs ``do not introduce any new divergences'' \cite[p.~431]{Ba3}; vacuum classes carry their vacuum counterterms \cite[p.~432]{Ba3}; and ``for an arbitrary graph $G$ we form a class $G_{\mathrm{ren}}$ containing this graph adding all graphs necessary to renormalize all divergent proper subgraphs. Of course localizations in these graphs should be in agreement'' \cite[p.~432]{Ba3}. For $G\in\mathcal C$, a pattern $\tau$ and a subgraph $S$ of $G$, say $\tau$ \emph{touches} $S$ if $\tau$ has an insertion on an internal scalar line of $S$ or a nonzero $\tilde B'$-order at a vertex of $S$; $S_\tau$ is the corresponding subgraph of $G_\tau$.

\begin{lemma}[the expanded class is a sum of renormalised classes]\label{sup3:lem:S144}
Let $\mathcal C$ be a per-label class (every member has the label multiset $\mathcal B$) at background $\tilde B$ on $\square$.
\begin{enumerate}[label=\textup{(\alph*)},nosep,leftmargin=*]
\item \emph{Untouched subgraphs.} Let $S$ be a divergent proper subgraph of $G\in\mathcal C$ and $G'\in\mathcal C$ a graph that differs from $G$ only inside $S$ (the counterterm graph $G/S$, or a leg-exchange or Ward partner inside $S$). The patterns of $G$ not touching $S$ correspond bijectively to the patterns of $G'$ trivial inside $S$, and $(G_\tau)'=(G')_\tau$. So inside $G_\tau$ the subgraph $S_\tau=S$ is renormalised by terms of the expansion of $\mathcal C$.
\item \emph{Touched subgraphs change type.} If $\tau$ touches $S$, then $S_\tau$ has $j\ge1$ more external vector legs than $S$ and $D(S_\tau)\ge D(S)+j$; $S_\tau$ is convergent if $\tau$ puts an $R$-term or an $F_2$-term inside $S$. The only partner sets a touched divergent $S_\tau$ can require are: (b1) $S$ of two-scalar-leg type with $j=1$: the partner with the two external scalar legs of $S_\tau$ exchanged; (b2) $S$ of type $1A$ with $j=1$, or a vacuum component with $j=2$: the Ward class of type $2A$; (b3) the four-leg graph \cite[(2.4)]{Ba3}, handled by the summation by parts of \cref{sup1:ssec:S4}. Touched $1A$ and $3A$ subgraphs need nothing.
\item \emph{Closure.} The per-label class $\mathcal C$ is closed under the exchange of the two external scalar legs of any two-scalar-leg or $\phi\phi A$ subgraph, the whole graph included (so a touched two-scalar-leg graph, now of type $\phi\phi A$ and degree $0$, finds its mirror partner in the expansion of $\mathcal C$), and the first (resp.\ second) $\tilde B'$-Taylor coefficient of the per-label $1A$ (resp.\ vacuum) class of an order is the per-label $2A$ class of that order with one (resp.\ two) external legs equal to $\tilde B'$, with the weights of the Taylor coefficients. Hence the partners of (b1)--(b2) are terms of the expansion of $\mathcal C$, and $\sum_{G,\tau}E(G_\tau)$ is a sum of per-label renormalised classes at background $B_0$, with $\tilde B'$ among the external fields.
\end{enumerate}
\end{lemma}

\begin{proof}
(a) Counterterm vertices \cite[(1.26)--(1.31)]{Ba3} and the vertices \cite[(1.6), (1.7), (1.10), (1.11)]{Ba3} do not depend on $\tilde B$. A pattern not touching $S$ consists of choices on the lines and vertices outside $S$, which are the same data for $G$ and $G'$; the expansion acts factor by factor (\cref{sup3:lem:S143}), so the two expansions agree outside $S$, and the labels of the new vertices are fixed by their own bonds and blocks, which also agree.

(b) The vertex degrees of \cite[p.~423]{Ba3} with the definition \cite[(2.2)]{Ba3}: a vertex of type \cite[(1.8)]{Ba3} or \cite[(1.10)]{Ba3} has degree $n(4-d)/2+n'$, so each added $\tilde B'$-leg raises the degree of the subgraph containing it by $1$, and an inserted vertex of \cref{sup3:lem:S142}(c) has $n=0$, $n'=j\ge1$ and adds $j$ to the degree of the line it sits on (two legs of dimension $-\frac{d-2}2$ each and one new line, net $j$); an $R$-term or a block factor inside $S$ makes $D(S_\tau)>0$ by \cite[Cor.~2.3]{Ba3} (\cref{sup2:lem:S7cor}). The list of divergent types with at most four external legs and their partner sets is \cite[pp.~429--431]{Ba3} (\cref{sup2:cor:S7class}): a two-scalar-leg subgraph of degree $-1$ touched once becomes $\phi\phi A$ of degree $0$, whose class is its symmetrisation; a $1A$ subgraph of degree $-1$ (resp.\ a vacuum component of degree $-2$) touched once (resp.\ twice) becomes $2A$ of degree $0$, whose class is the Ward class; a four-scalar-leg subgraph touched becomes \cite[(2.4)]{Ba3} with a vector leg, of degree $\ge1$, and is handled by \cref{sup1:lem:S41}; everything else touched has $D>0$.

(c) The two closure statements are the two lemmas on the class family in the fourth part of this appendix (\cref{sup4:sec}) (the exchange closure of per-label classes, which is a bijection of Wick pairings preserving orders and labels, and the identity of generating functions $\partial/\partial\tilde B'=\partial/\partial X$ for the Taylor coefficients of $\log$ of the Gaussian integral at background $B_0$, the vertices and lines depending on $\tilde B'$ and the other vector legs $X$ only through $B_0+\tilde B'+X$). Given them: for (b1), the exchange acts on the external connections of $S$ and the insertion on its internal lines, so $(S_\tau)^\pi=(S^\pi)_\tau$ with $S^\pi\subset G^\pi\in\mathcal C$; for (b2), the $2A$ Ward class at the order of $S_\tau$ is the set of all single (resp.\ double) $\tilde B'$-attachments, by line insertion or vertex expansion, to all graphs of $S$'s type and labels, which are exactly the touched terms of the expansion of $\mathcal C$. Together with (a) every divergent subgraph of every $G_\tau$ is renormalised, or completed to its Ward class, inside $\sum_{G,\tau}E(G_\tau)$, and the labels of the partners agree, so the sum is a sum of per-label renormalised classes.
\end{proof}

\begin{remark}\label{sup3:rem:S14b}
(a) The sentence ``which can be decomposed into a sum of renormalized classes'' of \cite[p.~433]{Ba3} is a closure property of the class family under differentiation in the background; it holds for per-label classes because these are complete at fixed labels (``if we refrain from taking the smallest classes'', \cite[p.~431]{Ba3}), and \cite[Prop.~1]{Ba3} for these classes is all that its consumers use, since they sum over all graphs of each order and over the localisations (\cite[p.~421]{Ba3}: ``Proposition I.3.1 is completely obvious if we sum over localizations''). (b) The regrouping and the renormalisation of \cite[pp.~434--445]{Ba3} are applied in the fourth part to per-label classes, so the classes produced here are the ones they take.
\end{remark}
